\section{Introduction}

A system of coupled harmonic oscillators decomposes into normal modes:
orthogonal coordinates that evolve independently, each with its own
characteristic frequency. The eigen-decomposition of an equity
correlation matrix has the same algebraic form. Its eigenvectors are
portfolios with (in-sample) uncorrelated returns, and its eigenvalues
are their variances. This paper asks whether the analogy can be made
\emph{predictive}. Do some normal modes of the market behave like
displaced oscillators that are pulled back to equilibrium, and can we
tell in advance which ones?

The question connects three literatures. Random matrix theory shows
that most of the eigenvalues of a large sample correlation matrix are
indistinguishable from noise, and the Marchenko--Pastur law gives a
principled edge between signal and noise
\citep{laloux1999,plerou2002,bunbouchaud2017}. Statistical arbitrage
trades the mean reversion of residual portfolios, typically residuals
of stocks against principal-component or ETF factors
\citep{avellaneda2010,gatev2006}, and short-term reversal profits are
known to vary strongly with market conditions
\citep{lehmann1990,nagel2012,khandani2011}. Machine-learning methods
have become a standard tool for extracting weak predictive signals from
many correlated characteristics \citep{gu2020}.

We make three contributions.

\emph{A physical prediction and a spectrum-wide test.} We take the
oscillator analogy to its overdamped, thermally driven limit, a
multivariate Ornstein--Uhlenbeck process. In that model the return
covariance and the stiffness matrix share eigenvectors and the
\emph{lowest}-variance modes are the fastest to revert. We then measure
out-of-sample reversion for every mode of the spectrum. The modes above
the Marchenko--Pastur edge, the conventional ``signal'', do not revert
(Sharpe \SRSignalLevel). Reversion is concentrated in the noise bulk
(Sharpe \SRMiddleLevel{} for modes 100--174, $t=\TMiddleLevel$), and
daily mode returns there have significantly negative autocorrelation.
The direction of the physical prediction holds, but its quantitative
form (half-life proportional to eigenvalue) does not.

\emph{Supervised prediction of mode quality.} Using features computed
only from each estimation window (eigenvalue separation, half-life,
autocorrelation, eigenvector stability, volatility, portfolio
concentration) plus market-regime variables, we train walk-forward
logistic, gradient-boosting and random-forest classifiers to predict
which candidate modes will revert profitably over the next month. The
signal is statistically significant but small (out-of-sample AUC
0.51--0.53).

\emph{A walk-forward strategy with regime integration.} We combine mode
predictions with filtered hidden-Markov regime probabilities in a
monthly re-estimated, market-neutral strategy, and evaluate returns,
turnover, drawdowns and robustness to costs, execution delay and
specification. Reversion is real before costs (gross Sharpe \GSREWB{}
when trading the full spectrum) but costly to harvest, with break-even
costs of only 1.7--2.5~bp per unit turnover. ML selection beats random selection
but not equal weighting. The largest improvement comes from regime
timing: reversion pays in stressed markets and loses in calm ones.

Throughout, we report negative and null results alongside positive
ones. The original specification in the accompanying code repository,
a daily-return reversal on the top signal modes, has a gross Sharpe of
\GSROrig{} on our data, and a central purpose of this paper is to
explain why.
